%!TEX root = ../main.tex
\TabAppendixResults
\renewcommand{\thefigure}{A\arabic{figure}}
\setcounter{figure}{0}

\section{Appendix}

\subsection{Dataset Statistics}

\begin{table}[H]
    \centering
    \begin{tabular}{lrrr}
        \toprule
        Dataset & Train   & Validation     & Test    \\
        \midrule
        PRIME   & 6{,}162 & 2{,}240 & 2{,}016 \\
        MAG     & 7{,}993 & 2{,}664 & 2{,}664 \\
        AMAZON  & 5{,}915 & 1{,}547 & 1{,}638 \\
        \bottomrule
    \end{tabular}
    \caption{Number of queries per dataset split for each STaRK graph.}
    \label{tab:traj_counts}
\end{table}

\newcommand{\TabStarkStatistics}[1][!ht]{%
    \begin{table}[#1]
        \centering
        \adjustbox{width=\columnwidth,center}{
            \begin{tabular}{lcccccc}
                \toprule
                Dataset 
                & \makecell{Entity\\types}
                & \makecell{Relation\\types}
                & \makecell{Average\\degree}
                & \makecell{Entities}
                & \makecell{Relations}
                & \makecell{Tokens} \\
                \midrule
                AMAZON & 4              & 5                & 18.2        & 1{,}035{,}542 & 9{,}443{,}802  & 592{,}067{,}882 \\
                MAG    & 4              & 4                & 43.5        & 1{,}872{,}968 & 39{,}802{,}116 & 212{,}602{,}571 \\
                PRIME  & 10             & 18               & 125.2       & 129{,}375     & 8{,}100{,}498  & 31{,}844{,}769  \\
                \bottomrule
            \end{tabular}
        }
        \caption{Statistics of the constructed semi-structured knowledge graphs used in STaRK.}
        \label{tab:stark-statistics}
    \end{table}%
}

\TabStarkStatistics

\subsection{Implementation Details}
\label{app:impl-details}

\retriever is run with $n{=}3$ parallel agents and a maximum trajectory length of $T_{\max}{=}20$. 

For the distilled variant, we use Qwen3-8B \cite{yang_qwen3_2025} (matching the model scale used by GraphFlow and Think-on-Graph with LLaMA3), trained via imitation on \retrieverShort teacher trajectories, while keeping the tool interface and exploration hyperparameters fixed.

Each tool call is executed with a bounded retrieval budget. Neighborhood exploration uses a fixed budget of $k{=}20$ neighbors per expansion. Global search returns up to $k$ nodes from the full graph. By default $k{=}5$, but the agent may override this value as a tool parameter.

Each agent outputs an ordered list of selected nodes. We aggregate these lists by ranking nodes first by the number of agents that selected them (vote count), and breaking ties by the earliest position at which the node appears in any agent’s list \cref{sec:parallel_exploration}. The aggregated ranking is truncated to the top 20 nodes to compute Recall@20; Hit@1 and MRR are computed on the same ranking.

\subsection{Additional Multi-Agent Analyses}
\label{app:budget-control}

\TabBudgetControlledComparison[h]

\noindent \Cref{tab:budget-controlled} compares a single longer trajectory with three shorter parallel trajectories under the same total step budget (30 steps). The three-agent configuration performs better on all three graphs, indicating that the gains from multi-agent inference are not explained solely by taking more total steps.

\subsection{Trajectory-Level Tool Usage}
\label{app:trajectory-tool-usage}

\TabTrajectoryToolMixing[h]

\noindent \Cref{tab:trajectory-tool-mixing} summarizes tool usage at the trajectory level. Because every trajectory begins with \textsc{Global Search}, runs either remain global-search-only or invoke both tools. On AMAZON, most trajectories, including successful ones, remain global-only, consistent with its text-heavy queries. On MAG and PRIME, by contrast, successful retrieval usually requires combining \textsc{Global Search} with neighborhood expansion.

\TabGlobalAfterFirstStep[h]

\noindent \Cref{tab:global-after-first-step} reports how often the agent returns to \textsc{Global Search} after the initial step. Re-anchoring is uncommon on AMAZON, but more frequent on MAG and especially PRIME, where retrieval more often requires shifting to a different region of the graph mid-trajectory. Together, \Cref{tab:trajectory-tool-mixing,tab:global-after-first-step} provide trajectory-level evidence for the two-tool design: AMAZON is often solved through lexical anchoring alone, whereas MAG and PRIME more often require combining global re-anchoring with local neighborhood expansion.

\subsection{Knowledge Graph Exploration Agent System Prompt}

The system prompt instructs the agent on the available node and relation types, the two tool signatures (\textsc{Global Search} and \textsc{Neighborhood Exploration}), output formatting, and the stopping criterion. The full prompt is included in the public repository of \retrieverShort at \url{https://github.com/mims-harvard/ark}.

\clearpage
%\FigArkAgentCli
